\subsection{Minors, Cofactors, and Rank}

\begin{definition} \label{def:minors-cofactors-and-adjugate} \hlt{Minors, Cofactors, and Adjugate}\\
Let $n\geq1$ and $A\in\operatorname{Mat}(n,n,\mathbb{F})$. Set $\det(\emptyset_{0\times0})\defeq1$.
\begin{enumerate}[label=(\roman*)]
\setlength{\itemsep}{0pt}
\item For $i,j\in\{1,\ldots,n\}$, the \hlt{minor matrix} $M_{ij}(A)$ is obtained by deleting row $i$ and column $j$ from $A$.\\
The corresponding \hlt{minor} is $\det(M_{ij}(A))$.
\item The $(i,j)$-\hlt{cofactor} is $C_{ij}(A)\defeq(-1)^{i+j}\det(M_{ij}(A))$.
\item The \hlt{cofactor matrix} is $C(A)\defeq[C_{ij}(A)]$, and the \hlt{adjugate} is $\operatorname{adj}(A)\defeq C(A)^t$.
\end{enumerate}
When $A$ is fixed, write $M_{ij}$ and $C_{ij}$.
\end{definition}

\begin{lemma} \label{lem:cofactor-expansion} \hlt{Cofactor Expansion}\\
Let $n\geq1$ and $A=[a_{ij}]\in\operatorname{Mat}(n,n,\mathbb{F})$.
\begin{enumerate}[label=(\roman*)]
\setlength{\itemsep}{0pt}
\item $(\forall i\in J_n)\,\det A=\sum\limits_{k=1}^{n}a_{ik}C_{ik}$.
\item $(\forall j\in J_n)\,\det A=\sum\limits_{k=1}^{n}a_{kj}C_{kj}$.
\end{enumerate}
\begin{proof}
\begin{enumerate}[label=(\roman*)]
\setlength{\itemsep}{0pt}
\item Fix $i\in J_n$. For $k\in J_n$, set $\mathfrak{S}_{i,k}\defeq\{\sigma\in\mathfrak{S}_n:\sigma(i)=k\}$ and $T_{i,k}\defeq\sum\limits_{\sigma\in\mathfrak{S}_{i,k}}\sign(\sigma)\prod\limits_{r=1}^{n}a_{r,\sigma(r)}$.\\
$(\forall k\in J_n)\,T_{i,k}=a_{ik}(-1)^{i+k}\det(M_{ik})=a_{ik}C_{ik}$ by the order-preserving deletion bijection $\mathfrak{S}_{i,k}\xrightarrow{\sim}\mathfrak{S}_{n-1}$, $i+k-2$ adjacent transpositions, the determinant and cofactor definitions, and Theorem \ref{thm:determinant-basic-laws}(ix).\\
$\det A=\sum\limits_{k\in J_n}T_{i,k}=\sum\limits_{k=1}^{n}a_{ik}C_{ik}$ by $\mathfrak{S}_n=\bigsqcup\limits_{k\in J_n}\mathfrak{S}_{i,k}$, Definition \ref{def:matrix-trace-and-determinant}(ii), Theorem \ref{thm:determinant-basic-laws}(ix), and the preceding identity.
\item Fix $j\in J_n$. $(\forall k\in J_n)\,[M_{jk}(A^t)=M_{kj}(A)^t\land C_{jk}(A^t)=C_{kj}(A)]$ by Definition \ref{def:minors-cofactors-and-adjugate} and Theorem \ref{thm:determinant-basic-laws}(ix).\\
$\det A=\det(A^t)=\sum\limits_{k=1}^{n}a_{kj}C_{jk}(A^t)=\sum\limits_{k=1}^{n}a_{kj}C_{kj}(A)$ by Theorem \ref{thm:determinant-basic-laws}(ix), part (i) for row $j$ of $A^t$, and the preceding identities.
\end{enumerate}
\hlt{Structural Commentary}: Permutation fibers prove the row formula; transposition transfers it to columns.\\
\hlt{Key Insight}: Deleting a fixed row and column leaves the corresponding minor sum.
\end{proof}
\end{lemma}

\begin{theorem} \label{thm:adjugate-identity-and-cramers-rule} \hlt{Adjugate Identity and Cramer's Rule}\\
Let $n\geq1$, $A\in\operatorname{Mat}(n,n,\mathbb{F})$, and $\mathbf{b}\in\mathbb{F}^n$. For $j\in\{1,\ldots,n\}$, let $A_j(\mathbf{b})$ be obtained from $A$ by replacing column $j$ by $\mathbf{b}$.
\begin{enumerate}[label=(\roman*)]
\setlength{\itemsep}{0pt}
\item $A\operatorname{adj}(A)=\operatorname{adj}(A)A=(\det A)I_n$.
\item $(\forall j\in\{1,\ldots,n\})\quad(\operatorname{adj}(A)\mathbf{b})_j=\det(A_j(\mathbf{b}))$.
\item If $\det A\neq0$, then $A^{-1}=(\det A)^{-1}\operatorname{adj}(A)$.
\item If $\det A\neq0$, then the unique solution of $A\mathbf{x}=\mathbf{b}$ satisfies $(\forall j\in\{1,\ldots,n\})\quad x_j=\frac{\det(A_j(\mathbf{b}))}{\det A}$.
\end{enumerate}
\begin{proof}
\begin{enumerate}[label=(\roman*)]
\setlength{\itemsep}{0pt}
\item For $i,j\in J_n$, $(A\operatorname{adj}(A))_{ij}=\sum\limits_{k=1}^{n}a_{ik}C_{jk}$ by Definition \ref{def:minors-cofactors-and-adjugate}.\\
The identity $\sum\limits_{k=1}^{n}a_{ik}C_{jk}=\delta_{ij}\det A$ holds by Lemma \ref{lem:cofactor-expansion}(i) and Theorem \ref{thm:determinant-basic-laws}(iv),(ix) with $A_{j,\cdot}\leftarrow A_{i,\cdot}$ when $i\neq j$.\\
Likewise, $(\operatorname{adj}(A)A)_{ij}=\sum\limits_{k=1}^{n}C_{ki}a_{kj}$ by Definition \ref{def:minors-cofactors-and-adjugate}.\\
The identity $\sum\limits_{k=1}^{n}C_{ki}a_{kj}=\delta_{ij}\det A$ holds by Lemma \ref{lem:cofactor-expansion}(ii) and Theorem \ref{thm:determinant-basic-laws}(iv) with $A_{\cdot,i}\leftarrow A_{\cdot,j}$ when $i\neq j$.\\
Hence $A\operatorname{adj}(A)=\operatorname{adj}(A)A=(\det A)I_n$.
\item Fix $j\in J_n$; $(\forall i\in J_n)\,C_{ij}(A_j(\mathbf{b}))=C_{ij}(A)$ by Definition \ref{def:minors-cofactors-and-adjugate}. Identity $(\operatorname{adj}(A)\mathbf{b})_j=\sum\limits_{i=1}^{n}C_{ij}(A)b_i$ holds by Definition \ref{def:minors-cofactors-and-adjugate}.\\
$\sum\limits_{i=1}^{n}C_{ij}(A)b_i=\sum\limits_{i=1}^{n}C_{ij}(A_j(\mathbf{b}))b_i=\det(A_j(\mathbf{b}))$ by the preceding cofactor identity and Lemma \ref{lem:cofactor-expansion}(ii).
\item If $\det A\neq0$, then $A((\det A)^{-1}\operatorname{adj}(A))=I_n=((\det A)^{-1}\operatorname{adj}(A))A$ by part (i).\\
Hence $A^{-1}=(\det A)^{-1}\operatorname{adj}(A)$ by the definition of a matrix inverse.
\item If $\det A\neq0$, then $\{\mathbf{x}\in\mathbb{F}^n:A\mathbf{x}=\mathbf{b}\}=\{A^{-1}\mathbf{b}\}$ by part (iii).\\
For $\mathbf{x}=A^{-1}\mathbf{b}$, $(\forall j\in J_n)\,x_j=(\det A)^{-1}(\operatorname{adj}(A)\mathbf{b})_j=\frac{\det(A_j(\mathbf{b}))}{\det A}$ by parts (ii)--(iii).
\end{enumerate}
\hlt{Structural Commentary}: Row and column replacements establish the adjugate identity before its inverse and solution consequences.\\
\hlt{Key Insight}: Off-diagonal cofactor sums create a repeated row or column.
\end{proof}
\end{theorem}

\begin{proposition} \label{prop:bordered-determinant-adjugate-form} \hlt{Adjugate Form of the Bordered Formula}\\
Let $n\geq1$, $A\in\operatorname{Mat}(n,n,\mathbb{F})$, $\mathbf{x},\mathbf{y}\in\mathbb{F}^n$, and $c\in\mathbb{F}$.\\
Then $\det\left[\begin{smallmatrix}A&\mathbf{x}\\ \mathbf{y}^t&c\end{smallmatrix}\right]=c\det(A)-\mathbf{y}^t\operatorname{adj}(A)\mathbf{x}$, without any invertibility assumption on $A$.
\begin{proof}
Set $D\defeq\left[\begin{smallmatrix}A&\mathbf{x}\\ \mathbf{y}^t&c\end{smallmatrix}\right]$. $(\forall j\in J_n)$, $M_{n+1,j}(D)=\left[A_{\cdot,1}\ \cdots\ A_{\cdot,j-1}\ A_{\cdot,j+1}\ \cdots\ A_{\cdot,n}\ \mathbf{x}\right]$.\\
$(\forall j\in J_n)$, $\det(M_{n+1,j}(D))=(-1)^{n-j}\det(A_j(\mathbf{x}))$ by $n-j$ column swaps and Theorem \ref{thm:determinant-basic-laws}(iii).\\
$(\forall j\in J_n)$, $C_{n+1,j}(D)=(-1)^{n+1+j+n-j}\det(A_j(\mathbf{x}))=-\det(A_j(\mathbf{x}))$ by Definition \ref{def:minors-cofactors-and-adjugate}.\\
$\det D=c\det A-\sum\limits_{j=1}^{n}y_j\det(A_j(\mathbf{x}))$ by Lemma \ref{lem:cofactor-expansion}(i) for row $n+1$ and the preceding cofactor identity.\\
$\det D=c\det A-\mathbf{y}^t\operatorname{adj}(A)\mathbf{x}$ by Theorem \ref{thm:adjugate-identity-and-cramers-rule}(ii) and the preceding identity.\\
\hlt{Structural Commentary}: Last-row expansion converts each border entry into a one-column replacement determinant.\\
\hlt{Key Insight}: The cofactor and column-move signs combine to one minus sign.
\end{proof}
\end{proposition}

\begin{theorem} \label{thm:cauchy-binet-minors} \hlt{Cauchy--Binet Formula}\\
Let $1\leq k\leq n$, $A\in\operatorname{Mat}(k,n,\mathbb{F})$, and $B\in\operatorname{Mat}(n,k,\mathbb{F})$. Let $\mathcal{J}_{k,n}$ be the set of $k$-element subsets of $\{1,\ldots,n\}$.\\
For $J=\{q_1<\cdots<q_k\}\in\mathcal{J}_{k,n}$, let $A^J$ consist of the columns of $A$ indexed by $J$, and let $B_J$ consist of the rows of $B$ indexed by $J$. Then $\det(AB)=\sum\limits_{J\in\mathcal{J}_{k,n}}\det(A^J)\det(B_J)$.
\begin{proof}
For $\mathbf{j}=(j_1,\ldots,j_k)\in J_n^k$, set $B_{\mathbf{j}}\defeq[b_{j_r,s}]_{r,s=1}^{k}$. Then $(\forall r\in J_k)\,(AB)_{r,\cdot}=\sum\limits_{j\in J_n}a_{rj}B_{j,\cdot}$.\\
Hence $\det(AB)=\sum\limits_{\mathbf{j}\in J_n^k}\left(\prod\limits_{r=1}^{k}a_{rj_r}\right)\det(B_{\mathbf{j}})$ by Theorem \ref{thm:determinant-basic-laws}(ii),(ix) and finite distributivity.\\
$[(\exists r,s\in J_k)\,(r\neq s\land j_r=j_s)]\Rightarrow\det(B_{\mathbf{j}})=0$ by Theorem \ref{thm:determinant-basic-laws}(iv),(ix).\\
Fix $J=\{q_1<\cdots<q_k\}\in\mathcal{J}_{k,n}$ and set $\mathbf{q}^{\tau}\defeq(q_{\tau^{-1}(r)})_{r\in J_k}$ for $\tau\in\mathfrak{S}_k$.\\
Then $\{\mathbf{j}\in J_n^k:\{j_r:r\in J_k\}=J\}=\{\mathbf{q}^{\tau}:\tau\in\mathfrak{S}_k\}$, and $(\forall\tau\in\mathfrak{S}_k)\,\det(B_{\mathbf{q}^{\tau}})=\sign(\tau)\det(B_J)$ by Theorem \ref{thm:determinant-basic-laws}(iii),(ix).\\
$\sum\limits_{\tau\in\mathfrak{S}_k}\sign(\tau)\left(\prod\limits_{r=1}^{k}a_{\tau(r),q_r}\right)\det(B_J)=\det(A^J)\det(B_J)$ by Definition \ref{def:matrix-trace-and-determinant}(ii).\\
Therefore $\det(AB)=\sum\limits_{J\in\mathcal{J}_{k,n}}\det(A^J)\det(B_J)$ by the preceding expansion, vanishing implication, and identities.\\
\hlt{Structural Commentary}: Row multilinearity expands over $J_n^k$; alternation restricts the sum to $\mathcal{J}_{k,n}$.\\
\hlt{Key Insight}: Each surviving subset contributes one column minor times one row minor.
\end{proof}
\end{theorem}

\begin{proposition} \label{prop:transpose-preserves-spectrum-and-geometric-multiplicity} \hlt{Transpose Preserves Spectrum and Geometric Multiplicity}\\
Let $n\geq1$, $A\in\operatorname{Mat}(n,n,\mathbb{F})$, and $\lambda\in\mathbb{F}$.
\begin{enumerate}[label=(\roman*)]
\setlength{\itemsep}{0pt}
\item $\dim(E_{\lambda}(A))=\dim(E_{\lambda}(A^t))$.
\item $\lambda\in\operatorname{spec}(A)$ iff $\lambda\in\operatorname{spec}(A^t)$.
\end{enumerate}
\begin{proof}
\begin{enumerate}[label=(\roman*)]
\setlength{\itemsep}{0pt}
\item[$(i)$] Since $(A-\lambda I_n)^t=A^t-\lambda I_n$, $\operatorname{rank}(A-\lambda I_n)=\operatorname{rank}(A^t-\lambda I_n)$ by Theorem \ref{thm:row-rank-equals-column-rank}.\\
Hence $\dim(E_{\lambda}(A))=n-\operatorname{rank}(A-\lambda I_n)=n-\operatorname{rank}(A^t-\lambda I_n)=\dim(E_{\lambda}(A^t))$ by Corollary \ref{cor:matrix-rank-nullity}(i).
\item[$(ii)$] ($\Rightarrow$) If $\lambda\in\operatorname{spec}(A)$, then $\dim(E_{\lambda}(A))>0$ by the matrix-spectrum definition.\\
Hence $\dim(E_{\lambda}(A^t))>0$ by part (i), so $\lambda\in\operatorname{spec}(A^t)$ by the matrix-spectrum definition.
\item[$(ii)$] ($\Leftarrow$) If $\lambda\in\operatorname{spec}(A^t)$, then $\dim(E_{\lambda}(A^t))>0$ by the matrix-spectrum definition.\\
Hence $\dim(E_{\lambda}(A))>0$ by part (i), so $\lambda\in\operatorname{spec}(A)$ by the matrix-spectrum definition.
\end{enumerate}
\hlt{Structural Commentary}: Transpose-invariant rank fixes shifted-kernel dimensions before the two spectral directions.\\
\hlt{Key Insight}: Spectrum and geometric multiplicity are nontriviality and dimension of shifted kernels.
\end{proof}
\end{proposition}

\begin{corollary} \label{cor:matrix-nonsingularity-rank-determinant-spectrum-dictionary} \hlt{Nonsingularity by Rank, Determinant, and Spectrum}\\
Let $n\geq1$ and $A\in\operatorname{Mat}(n,n,\mathbb{F})$. The following are equivalent:
\begin{enumerate}[label=(\roman*)]
\setlength{\itemsep}{0pt}
\item $A$ is invertible;
\item $\operatorname{rank}(A)=n$;
\item $\det A\neq0$;
\item $0\notin\operatorname{spec}(A)$.
\end{enumerate}
\begin{proof}
\begin{enumerate}[label=(\roman*)]
\setlength{\itemsep}{0pt}
\item[$(i)$]$\Rightarrow (ii)$ If $A$ is invertible, then $\operatorname{rank}(A)=n$ by Corollary \ref{cor:square-matrix-invertible-iff-full-rank}. \item[$(ii)$]$\Rightarrow (iii)$ If $\operatorname{rank}(A)=n$, then $A$ is invertible by Corollary \ref{cor:square-matrix-invertible-iff-full-rank}. Hence $\det A\neq0$ by Theorem \ref{thm:determinant-basic-laws}(vii). \item[$(iii)$]$\Rightarrow (iv)$ If $\det A\neq0$, then $A$ is invertible by Theorem \ref{thm:determinant-basic-laws}(vii). Hence $0\notin\operatorname{spec}(A)$ by Corollary \ref{cor:matrix-invertible-iff-zero-not-spectrum}. \item[$(iv)$]$\Rightarrow (i)$ If $0\notin\operatorname{spec}(A)$, then $A$ is invertible by Corollary \ref{cor:matrix-invertible-iff-zero-not-spectrum}.
\end{enumerate}
\hlt{Structural Commentary}: The cycle passes through full rank and determinant before returning through zero-spectrum exclusion.\\
\hlt{Key Insight}: Invertibility has equivalent range, determinant, and shifted-kernel tests.
\end{proof}
\end{corollary}

\begin{theorem} \label{thm:rank-product-and-concatenation-bounds} \hlt{Rank Bounds for Products and Concatenations}\\
Let $A\in\operatorname{Mat}(m,k,\mathbb{F})$, $B\in\operatorname{Mat}(k,n,\mathbb{F})$, and $C\in\operatorname{Mat}(m,p,\mathbb{F})$. Set $D\defeq\left[\begin{smallmatrix}A&C\end{smallmatrix}\right]$.
\begin{enumerate}[label=(\roman*)]
\setlength{\itemsep}{0pt}
\item $\operatorname{rank}(AB)\leq\min(\operatorname{rank}(A),\operatorname{rank}(B))$.
\item $\operatorname{rank}(AB)=\operatorname{rank}(A)\Longleftrightarrow\operatorname{col}(AB)=\operatorname{col}(A)$.
\item $\max(\operatorname{rank}(A),\operatorname{rank}(C))\leq\operatorname{rank}(D)$.
\item $\operatorname{rank}(D)=\operatorname{rank}(A)\Longleftrightarrow\operatorname{col}(C)\subseteq\operatorname{col}(A)$.
\end{enumerate}
\begin{proof}
Column-space inclusion is rank-monotone, with equality iff the spaces coincide, by the rank definition and Proposition \ref{prop:dimension-bounds-and-basis-criteria}.\\
For matrices $X,Y$ of compatible dimensions, $\operatorname{col}(XY)\subseteq\operatorname{col}(X)$ by Propositions \ref{prop:matrix-product-by-columns} and \ref{prop:matrix-vector-product-column-combination}.\\
The identities $(XY)^t=Y^tX^t$ and $\operatorname{rank}(X^t)=\operatorname{rank}(X)$ hold by Proposition \ref{prop:transpose-reverses-products} and Theorem \ref{thm:row-rank-equals-column-rank}, respectively.\\
The identity $\operatorname{col}(D)=\operatorname{col}(A)+\operatorname{col}(C)$ holds by the definition of $D$ and the column-space definition.
\begin{enumerate}[label=(\roman*)]
\setlength{\itemsep}{0pt}
\item[$(i)$] $\operatorname{rank}(AB)\leq\operatorname{rank}(A)$ by the product-column inclusion and the rank criterion.\\
$\operatorname{rank}(AB)\leq\operatorname{rank}(B)$ by the transpose identities, the product-column inclusion for $B^tA^t$, and the rank criterion.
\item[$(ii)$] ($\Rightarrow$) $\operatorname{col}(AB)=\operatorname{col}(A)$ by $\operatorname{rank}(AB)=\operatorname{rank}(A)$, the product-column inclusion, and the rank criterion.
\item[$(ii)$] ($\Leftarrow$) $\operatorname{rank}(AB)=\operatorname{rank}(A)$ by $\operatorname{col}(AB)=\operatorname{col}(A)$ and the rank definition.
\item[$(iii)$] $\max(\operatorname{rank}(A),\operatorname{rank}(C))\leq\operatorname{rank}(D)$ by $\operatorname{col}(A)\cup\operatorname{col}(C)\subseteq\operatorname{col}(D)$, the concatenation identity, and the rank criterion.
\item[$(iv)$] ($\Rightarrow$) $\operatorname{col}(C)\subseteq\operatorname{col}(D)=\operatorname{col}(A)$ by the concatenation identity, $\operatorname{rank}(D)=\operatorname{rank}(A)$, and the rank criterion.
\item[$(iv)$] ($\Leftarrow$) $\operatorname{col}(C)\subseteq\operatorname{col}(A)\Rightarrow\operatorname{col}(D)=\operatorname{col}(A)$ by the concatenation identity.\\
Hence $\operatorname{rank}(D)=\operatorname{rank}(A)$ by the rank definition.
\end{enumerate}
\hlt{Structural Commentary}: Product and concatenation identities reduce every claim to nested column spaces.\\
\hlt{Key Insight}: Equal dimensions force equality after column-space inclusion is established.
\end{proof}
\end{theorem}

\begin{theorem} \label{thm:sylvester-rank-inequality} \hlt{Sylvester Rank Inequality}\\
Let $A\in\operatorname{Mat}(m,k,\mathbb{F})$ and $B\in\operatorname{Mat}(k,n,\mathbb{F})$. Then $\operatorname{rank}(A)+\operatorname{rank}(B)-k\leq\operatorname{rank}(AB)$.
\begin{proof}
Set $R\defeq\operatorname{col}(B)$ and define $S:R\to\mathbb{F}^m$ by $S(\mathbf{x})\defeq A\mathbf{x}$.\\
The identities $\operatorname{range}(S)=\operatorname{col}(AB)$ and $\ker(S)=R\cap\ker(A)$ hold by Proposition \ref{prop:matrix-product-by-columns} and the definitions of $R$ and $S$.\\
$\operatorname{rank}(AB)=\dim(R)-\dim(R\cap\ker(A))\geq\operatorname{rank}(B)-\dim(\ker(A))$ by Theorem \ref{thm:rank-nullity}(ii), the rank definition, and Proposition \ref{prop:dimension-bounds-and-basis-criteria}(i).\\
$\operatorname{rank}(AB)\geq\operatorname{rank}(B)-(k-\operatorname{rank}(A))=\operatorname{rank}(A)+\operatorname{rank}(B)-k$ by the preceding bound and Corollary \ref{cor:matrix-rank-nullity}(i).\\
\hlt{Structural Commentary}: Rank--nullity restricts $A$ to $\operatorname{col}(B)$ and isolates the lost middle-space directions.\\
\hlt{Key Insight}: Only $\operatorname{col}(B)\cap\ker(A)$ can disappear under multiplication.
\end{proof}
\end{theorem}

\begin{proposition} \label{prop:sylvester-law-of-nullity} \hlt{Sylvester's Law of Nullity}\\
For a matrix $X$, write $\operatorname{nullity}(X)\defeq\dim(\ker(X))$.
\begin{enumerate}[label=(\roman*)]
\setlength{\itemsep}{0pt}
\item If $A\in\operatorname{Mat}(m,k,\mathbb{F})$ and $B\in\operatorname{Mat}(k,n,\mathbb{F})$, then $\operatorname{nullity}(AB)\leq\operatorname{nullity}(A)+\operatorname{nullity}(B)$.
\item If $A,B\in\operatorname{Mat}(n,n,\mathbb{F})$, then $\max(\operatorname{nullity}(A),\operatorname{nullity}(B))\leq\operatorname{nullity}(AB)$.
\end{enumerate}
\begin{proof}
\begin{enumerate}[label=(\roman*)]
\setlength{\itemsep}{0pt}
\item $\operatorname{nullity}(AB)\leq n-\operatorname{rank}(A)-\operatorname{rank}(B)+k=\operatorname{nullity}(A)+\operatorname{nullity}(B)$ by Theorem \ref{thm:sylvester-rank-inequality} and Corollary \ref{cor:matrix-rank-nullity}(i) for $AB$, $A$, and $B$.
\item $\operatorname{nullity}(AB)\geq\max(n-\operatorname{rank}(A),n-\operatorname{rank}(B))=\max(\operatorname{nullity}(A),\operatorname{nullity}(B))$ by Theorem \ref{thm:rank-product-and-concatenation-bounds}(i) and Corollary \ref{cor:matrix-rank-nullity}(i).
\end{enumerate}
\hlt{Structural Commentary}: Rank--nullity reverses the lower and upper product-rank bounds into nullity bounds.\\
\hlt{Key Insight}: Subtracting rank from a fixed domain dimension reverses inequalities.
\end{proof}
\end{proposition}

\begin{theorem} \label{thm:frobenius-rank-inequality} \hlt{Frobenius Rank Inequality}\\
Let $A\in\operatorname{Mat}(m,k,\mathbb{F})$, $B\in\operatorname{Mat}(k,p,\mathbb{F})$, and $C\in\operatorname{Mat}(p,n,\mathbb{F})$. Then $\operatorname{rank}(AB)+\operatorname{rank}(BC)\leq\operatorname{rank}(B)+\operatorname{rank}(ABC)$.
\begin{proof}
Set $R\defeq\operatorname{col}(B)$ and $S\defeq\operatorname{col}(BC)$, and define $T_A:\mathbb{F}^k\to\mathbb{F}^m$ by $T_A(\mathbf{x})\defeq A\mathbf{x}$.\\
The identities $S\subseteq R$, $T_A(R)=\operatorname{col}(AB)$, and $T_A(S)=\operatorname{col}(ABC)$ hold by Proposition \ref{prop:matrix-product-by-columns}.\\
$(\forall U\in\{R,S\})\,\dim(T_A(U))=\dim(U)-\dim(U\cap\ker(A))$ by Theorem \ref{thm:rank-nullity}(ii) for $T_A|_U$ and the definition of restriction.\\
The Frobenius inequality holds by $S\cap\ker(A)\subseteq R\cap\ker(A)$, Proposition \ref{prop:dimension-bounds-and-basis-criteria}(i), the preceding identities, and the rank definition.\\
\hlt{Structural Commentary}: Rank--nullity compares $A$ on the nested column spaces of $BC$ and $B$.\\
\hlt{Key Insight}: The smaller space cannot meet $\ker(A)$ in more dimensions.
\end{proof}
\end{theorem}

Theorem \ref{thm:frobenius-rank-inequality} contains both two-factor bounds as degenerations: a zero outer factor recovers Theorem \ref{thm:rank-product-and-concatenation-bounds}(i), while an identity middle factor recovers Theorem \ref{thm:sylvester-rank-inequality}.

\begin{theorem} \label{thm:rank-triangle-inequality} \hlt{Rank Triangle Inequality}\\
Let $A,B\in\operatorname{Mat}(m,n,\mathbb{F})$.
\begin{enumerate}[label=(\roman*)]
\setlength{\itemsep}{0pt}
\item $\operatorname{rank}(A+B)\leq\operatorname{rank}(A)+\operatorname{rank}(B)$.
\item $\abs{\operatorname{rank}(A)-\operatorname{rank}(B)}\leq\operatorname{rank}(A+B)$.
\end{enumerate}
\begin{proof}
\begin{enumerate}[label=(\roman*)]
\setlength{\itemsep}{0pt}
\item $\operatorname{col}(A+B)\subseteq\operatorname{col}(A)+\operatorname{col}(B)$ by adding corresponding columns.\\
$\operatorname{rank}(A+B)\leq\operatorname{rank}(A)+\operatorname{rank}(B)$ by Proposition \ref{prop:dimension-bounds-and-basis-criteria}(i), Theorem \ref{thm:dimension-formula-for-sum}, and the rank definition.
\item $\abs{\operatorname{rank}(A)-\operatorname{rank}(B)}\leq\operatorname{rank}(A+B)$ by part (i) with $A=(A+B)+(-B)$,
and $B=(A+B)+(-A)$, $(\forall X\in\{A,B\})\,\operatorname{col}(-X)=\operatorname{col}(X)$, and the rank and absolute-value definitions.
\end{enumerate}
\hlt{Structural Commentary}: The column-space sum proves the upper bound; rearranging each summand proves both reverse bounds.\\
\hlt{Key Insight}: Rewriting either summand converts the upper bound into a lower bound.
\end{proof}
\end{theorem}
